%---------------------------------------------------------------------
% Title page and one orientation page.
%---------------------------------------------------------------------
\thispagestyle{empty}

% Band depth measured from the rendered page, with headroom. Nastaliq at
% Scale=1.45 sets far taller lines than the nominal point sizes suggest.
\begin{tikzpicture}[remember picture, overlay]
  \fill[bmblued] (current page.north west) rectangle
        ([yshift=-12.3cm]current page.north east);
  \fill[bmaccent] ([yshift=-12.3cm]current page.north west) rectangle
        ([yshift=-12.5cm]current page.north east);
\end{tikzpicture}

\vspace*{0.9cm}
{\color{white}
  {\large علامہ اقبال اوپن یونیورسٹی، اسلام آباد}\par
  \vspace{3pt}
  {شعبۂ ابلاغِ عامہ \ \textbullet\ \ بی اے}\par
  \vspace{20pt}
  {\Huge\bfseries خبر نگاری}\par
  \vspace{6pt}
  {\large\bfseries \textenglish{Reporting}}\par
  \vspace{12pt}
  {\LARGE\bfseries متوقع سوالات اور ان کے جوابات}\par
  \vspace{12pt}
  {\large کورس کوڈ \textenglish{431}}\par
  \vspace{4pt}
  {چودہ سوالات \ \textbullet\ \ دو حقیقی پرچوں سے ماخوذ\par}
}

\vspace{1.0cm}

\begin{tcolorbox}[enhanced, colback=bmtint, colframe=bmblue, boxrule=0.6pt,
                  arc=3pt, left=12pt, right=12pt, top=10pt, bottom=10pt]
\textbf{\color{bmblue}سوالات کہاں سے لیے گئے ہیں؟}\ \
تمام سوالات اے آئی او یو کے \textbf{دو حقیقی گزشتہ پرچوں} سے لفظ بہ لفظ نقل
کیے گئے ہیں:

\smallskip
\begin{itemize}\setlength{\itemsep}{0pt}\setlength{\topsep}{2pt}
  \item بی اے، خبر نگاری (\textenglish{431})، سمسٹر \textbf{بہار
        \textenglish{2013}} --- آٹھ سوالات
  \item بی اے، خبر نگاری (\textenglish{431})، سمسٹر \textbf{خزاں
        \textenglish{2013}} --- آٹھ سوالات
\end{itemize}

\smallskip
دونوں پرچوں میں دہرائے گئے سوالات کو یکجا کر دیا گیا ہے --- سولہ میں سے چودہ
الگ موضوعات بنے، اور یہی دہراؤ بتاتا ہے کہ کون سا موضوع سب سے اہم ہے۔

\medskip
\textbf{\color{bmblue}ایک صاف بات۔}\ \
اس کورس کی سرکاری درسی کتاب دستیاب نہ ہو سکی، اس لیے یہ کتابچہ کتاب کے اختتامی
سوالات پر نہیں بلکہ \emph{گزشتہ پرچوں} پر مبنی ہے۔ جوابات خبر نگاری کے مسلمہ
اصولوں اور صحافت کی معلوم تاریخ کے مطابق ہیں۔ جہاں کتاب کی مخصوص تفصیل درکار
ہو، اپنی درسی کتاب سے ملا لیجیے۔

\medskip
\textbf{\color{bmblue}حیثیت۔}\ \
ذاتی مطالعے کے لیے۔ یہ اے آئی او یو کی مطبوعہ کتاب نہیں۔ اسے من و عن نقل کر کے
اسائنمنٹ میں جمع کرانا منع ہے۔
\end{tcolorbox}

\clearpage

%---------------------------------------------------------------------
\section*{پرچے کا خاکہ}
\markboth{پرچے کا خاکہ}{}

دونوں پرچوں کی ساخت بالکل یکساں ہے:

\medskip
\begin{center}
\begin{tabular}{@{}p{4.6cm} p{9.8cm}@{}}
\toprule
کل نمبر & \textenglish{100} \\
کامیابی کے نمبر & \textenglish{40} \\
وقت & تین گھنٹے \\
سوالات دیے جاتے ہیں & آٹھ \\
کتنے حل کرنے ہیں & \textbf{کوئی سے پانچ} \\
فی سوال نمبر & \textenglish{20} --- تمام سوالوں کے نمبر برابر ہیں \\
\bottomrule
\end{tabular}
\end{center}

\medskip
\begin{nukta}[خبر نگاری کے پرچے کی خاص بات]
اس مضمون میں محض تعریف لکھ دینا کافی نہیں --- ممتحن \emph{عملی صحافتی سمجھ}
دیکھتا ہے۔ اس لیے ہر جواب میں:

\smallskip
\textbf{(۱)} تعریف، \textbf{(۲)} اصول یا اقسام \emph{عنوانات کے ساتھ}،
\textbf{(۳)} کم از کم ایک \textbf{عملی مثال} --- کوئی حقیقی اخبار، ادارہ یا
صورتِ حال، اور \textbf{(۴)} مختصر خلاصہ۔

\smallskip
جہاں ممکن ہو، خبر کا سادہ ڈھانچہ (سرخی، انٹرو، باڈی) یا الٹے اہرام کا خاکہ بنا
دینا اضافی نمبر دلاتا ہے۔
\end{nukta}

\medskip
\noindent
\textbf{\color{bmblue}کون سے موضوعات دونوں پرچوں میں آئے}

\medskip
\begin{center}
\begin{tabular}{@{}p{6.6cm} p{2.6cm} p{5.2cm}@{}}
\toprule
\textbf{موضوع} & \textbf{کتنے پرچوں میں} & \textbf{اس کتابچے میں سوال} \\
\midrule
رپورٹر --- فرائض، اوصاف اور اقسام کا فرق & دونوں & \textenglish{6}، \textenglish{7} \\
خبر کا مفہوم اور کردار                    & دونوں & \textenglish{1} \\
مکتوب نگاری                                & دونوں & \textenglish{12} \\
خبر رساں ادارے                             & دونوں & \textenglish{13} \\
ترقی یافتہ / تجزیاتی خبر نگاری             & دونوں & \textenglish{10}، \textenglish{11} \\
خبر کی اقدار اور اقسام                     & بہار  & \textenglish{2}، \textenglish{3} \\
خبر نگاری کے اصول و اسلوب                  & بہار  & \textenglish{4} \\
بیٹ رپورٹنگ                                 & بہار  & \textenglish{8} \\
خبروں کے حصول کے ذرائع                     & خزاں  & \textenglish{5} \\
علاقائی نامہ نگاری                          & خزاں  & \textenglish{9} \\
قوانینِ صحافت                               & خزاں  & \textenglish{14} \\
\bottomrule
\end{tabular}
\end{center}

\medskip
\noindent
پہلی پانچ سطریں دونوں پرچوں میں آئی ہیں۔ اگر وقت بہت کم ہو تو انہی سے شروع
کیجیے --- پانچ میں سے تین سوال عموماً انہی سے بن جاتے ہیں۔
